% Generated draft: requires mathematical and source review.
\documentclass[11pt]{article}
\usepackage[margin=25mm]{geometry}
\usepackage{fontspec}
\setmainfont{Latin Modern Roman}
\usepackage{amsmath,amssymb,mathtools,unicode-math}
\setmathfont{Latin Modern Math}
\usepackage{enumitem,needspace,xurl,hyperref}
\setlength{\parindent}{0pt}
\setlength{\parskip}{5pt}
\setlength{\emergencystretch}{3em}
\newcommand{\N}{\mathbb{N}}
\newcommand{\Z}{\mathbb{Z}}
\newcommand{\Q}{\mathbb{Q}}
\newcommand{\R}{\mathbb{R}}
\newcommand{\C}{\mathbb{C}}
\newcommand{\F}{\mathbb{F}}
\newcommand{\A}{\mathbb{A}}
\newcommand{\PP}{\mathbb{P}}
\newcommand{\E}{\mathbb{E}}
\newcommand{\eps}{\varepsilon}
\newcommand{\dd}{\,\mathrm{d}}
\begin{document}
\section*{UnsolvedMath ID: 30006417}
\textbf{Source problem code:} catalogue entry
\subsection{Definitions and mathematical statement}
Let \(M\) be a smooth closed three-manifold diffeomorphic to \(S^3\), and let \(h\) be a smooth Riemannian metric on \(M\). Here closed means compact without boundary. For each integer \(g\geq 0\), let \(\Sigma_g\) denote a smooth closed connected orientable surface of genus \(g\). A smooth embedding \(f:\Sigma_g\to M\) is minimal if its mean curvature vector with respect to \(h\) vanishes identically. Define
\[
\mathcal M_g^{\mathrm{emb}}(M,h)
=\bigl\{f(\Sigma_g): f:\Sigma_g\hookrightarrow M
\text{ is a smooth minimal embedding}\bigr\}.
\]
Elements are counted as distinct subsets of \(M\), so reparametrizations of one surface do not give additional surfaces. Thus \(\Sigma_0\cong S^2\) and \(\Sigma_1\cong S^1\times S^1\).

Under the embedded-surface interpretation, the two principal questions attributed in the supplied excerpt to Yau (1982), for spheres, and White (1989), for tori, are whether
\[
\forall (M,h)\text{ as above},\qquad
\#\mathcal M_0^{\mathrm{emb}}(M,h)\geq 4
\quad\text{and}\quad
\#\mathcal M_1^{\mathrm{emb}}(M,h)\geq 5.
\]
There is no curvature or genericity assumption in this unrestricted question. The quantifier includes metrics outside the positive-Ricci-curvature cases mentioned in the source.

The supplied record also asks for multiple minimal surfaces of every genus. If ``multiple'' means at least two distinct embedded images, its weakest numerical formulation is
\[
\forall (M,h)\text{ as above}\;\forall g\in\Z_{\geq0},
\qquad \#\mathcal M_g^{\mathrm{emb}}(M,h)\geq 2.
\]
This is only a provisional interpretation: the excerpt anticipates a number related to the topology of the space of genus-\(g\) Lawson surfaces, but gives neither a numerical bound nor a definition of that space or the relevant topological invariant. Moreover, the supplied wording does not explicitly specify embedded versus immersed surfaces or its convention for distinctness. The displayed formulations make the embedded interpretation explicit; these conventions require confirmation before adoption as definitive statements.

The supplied dated curation labels the problem partially solved and reports the sphere and torus bounds in positive Ricci curvature, with the sphere bound also reported for bumpy metrics. The excerpt attributes the positive-Ricci sphere result to Wang--Zhou (2023) and the torus result to work of Yangyang Li with Adrian Chu (2024). It also mentions techniques for higher genus, particularly genus \(2\), in joint work with Adrian Chun-Pong Chu and Zhihan Wang, without supplying a precise higher-genus theorem. These are supplied contextual claims, not independently verified results of this draft.
\subsection{Short English statement}
The question asks whether every smooth Riemannian three-sphere contains at least four distinct minimal spheres and five distinct minimal tori, with no curvature restriction. It further asks for multiple minimal surfaces of every genus. An embedded-surface formulation is given explicitly, but the source leaves embeddedness, distinctness, and the intended higher-genus lower bounds unresolved. The supplied partially solved status and special-case results require independent review.
\subsection{Sources}
\begin{itemize}
\item Supplied local UnsolvedMath registry record, permanent ID \(30006417\), title ``Multiple Minimal Surfaces in Riemannian Three-Spheres,'' source problem code \texttt{catalogue entry}, in the Oberwolfach Reports Open Problems collection. The registry supplies the status \texttt{partially\_solved} and the editorial label \texttt{corrected\_verified}; these labels do not constitute independent mathematical verification.
\item Supplied report citation: \emph{Differentialgeometrie im Großen (2026)}, Oberwolfach Reports, DOI \texttt{10.4171/owr/2025/30}, pp.\ 1591--1646. \url{https://doi.org/10.4171/owr/2025/30}. The literature-source list calls this ``Oberwolfach Report 30/2025.'' The supplied citation's year and the report identifier differ and have not been reconciled.
\item Supplied excerpt attribution: Yangyang Li, ``Constructing minimal surfaces of prescribed genus in closed Riemannian 3-spheres with positive Ricci curvature,'' joint work with Adrian Chun-Pong Chu and Zhihan Wang. The excerpt attributes the sphere conjecture to Yau (1982), reference [4], and the torus conjecture to White (1989), reference [5]. It attributes the construction of arbitrary-genus Lawson surfaces in \(S^3\) to Lawson (1970), reference [6]. Full bibliographic entries for these numbered references are not supplied.
\item Supplied literature source: ``Existence of four minimal spheres in S3 with a bumpy metric.'' \url{https://arxiv.org/abs/2305.08755}. The excerpt separately cites Wang--Zhou (2023) as reference [7]; its complete bibliography is not supplied.
\item Supplied literature source: ``Existence of 5 minimal tori in 3-spheres of positive Ricci curvature.'' \url{https://arxiv.org/abs/2409.09315}. The excerpt describes a 2024 result of Adrian Chu and Yangyang Li as reference [8], and higher-genus joint work with Adrian Chu and Zhihan Wang as reference [9]; complete entries are not supplied.
\item Literature assessment date: \texttt{2026-08-22}. This is supplied dated curation, not an independently verified assessment. It reports the positive-Ricci cases and the bumpy-metric sphere case as proved, and the unrestricted claim as open. The background explicitly requests independent verification of time-sensitive status and cited claims. No linked source was visited for this draft.
\end{itemize}
\end{document}
